\documentclass[11pt, a4paper]{article}

% --- Pakete ---
\usepackage[english, ngerman]{babel}
\usepackage[utf8]{inputenc}
\usepackage[T1]{fontenc}
\usepackage{lmodern}
\usepackage[left=2.2cm, right=2.2cm, top=2.2cm, bottom=2.2cm]{geometry}
\usepackage[autostyle=true,german=quotes]{csquotes}
\usepackage{graphicx}
\usepackage{xcolor}
\usepackage{enumitem}
\usepackage{fancyhdr}
\usepackage{titlesec}
\usepackage{microtype}
\usepackage{pdflscape}
\usepackage{float}
\usepackage[font=small, labelfont={bf,color=hkared}, skip=6pt]{caption}
\usepackage{tikz}
\usetikzlibrary{arrows.meta, calc, shapes.symbols}
\usepackage[document]{ragged2e}
\usepackage{hyperref}
\usepackage{xurl}

\emergencystretch=0.3\textwidth
\hyphenpenalty=1000
\exhyphenpenalty=50

% --- Farben (HKA Corporate) ---
\definecolor{hkared}{RGB}{198, 42, 18}
\definecolor{hkagray}{RGB}{100, 100, 100}

% --- TikZ-Stile (angelehnt an die Themenlandkarten-Beispiele der Vorlesung) ---
% Jede Kante ist gerichtet und beschriftet
\tikzset{
    begriff/.style = {draw=black, line width=1.3pt, fill=white, align=center,
                      font=\small\bfseries, minimum width=2.8cm, minimum height=0.95cm},
    pfeil/.style   = {-{Stealth[length=2.4mm, width=1.8mm]}, thin, draw=black},
    txt/.style     = {font=\scriptsize, fill=white, inner sep=1.5pt, align=center},
    kbox/.style    = {begriff, fill=#1!12, text width=2.8cm, minimum width=0cm,
                      font=\footnotesize\bfseries, minimum height=0.9cm},
}

% --- Farben der Bereiche (ausführliche Karte) ---
\definecolor{clA}{RGB}{31, 95, 155}   % Gesellschaft
\definecolor{clB}{RGB}{198, 42, 18}   % Menschen
\definecolor{clC}{RGB}{46, 125, 50}   % Waren
\definecolor{clD}{RGB}{200, 120, 0}   % Hindernisse

% --- Hyperref Setup ---
\hypersetup{
    colorlinks  = true,
    linkcolor   = hkared,
    urlcolor    = hkared,
    citecolor   = hkared,
    pdftitle    = {Zeitgeist und Themenkarte -- Tauschhandel (Kurzfassung)},
    pdfauthor   = {Erlind Sejdiu},
    pdfsubject  = {Game Design -- Game Concept Document, Kapitel 1},
    pdfcreator  = {pdflatex},
    pdfproducer = {LaTeX with hyperref}
}

% --- Section Styling ---
\titleformat{\section}{\large\bfseries\color{hkared}}{\thesection}{0.8em}{}
\titleformat{\subsection}{\normalsize\bfseries\color{hkared}}{\thesubsection}{0.8em}{}
\titlespacing*{\section}{0pt}{0.85em}{0.25em}
\titlespacing*{\subsection}{0pt}{0.5em}{0.15em}

% --- Header/Footer ---
\pagestyle{fancy}
\fancyhf{}
\fancyhead[L]{\small\color{hkagray}Labor Game Design}
\fancyhead[R]{\small\color{hkagray}Master Informatik WS~2026/27}
\fancyfoot[C]{\small\color{hkagray}\thepage}
\renewcommand{\headrulewidth}{0.4pt}

% Listenabstände straffen
\setlist{noitemsep, topsep=2pt, parsep=1pt, partopsep=0pt}

% Abstand zwischen Absätzen
\setlength{\parskip}{0.6em}

% =============================================================================
\begin{document}

% --- Titelseite ---
\hypersetup{pageanchor=false}
\begin{titlepage}
    \centering
    \includegraphics[width=0.35\textwidth]{HKA.jpg}

    {\normalsize\color{hkagray} University of Applied Sciences}\\[0.125cm]
    {\small\color{hkagray} Fakultät für Informatik und Wirtschaftsinformatik}

    \vspace{.5cm}
    {\color{hkared}\rule{\textwidth}{1.5pt}}

    \vspace{1.5cm}
    {\fontsize{24}{30}\selectfont\bfseries\color{hkared} Zeitgeist und Themenkarte -- Tauschhandel}\\[0.6cm]
    {\large\color{hkagray} Game Concept Document -- Kapitel 1 (Kurzfassung)}
    \vspace{.75cm}

    \vspace{.5cm}
    {\large Labor Game Design}\\[0.3cm]
    {\large Wintersemester 2026/27}

    \vspace{.5cm}
    {\color{hkared}\rule{0.6\textwidth}{0.5pt}}
    \vspace{1.5cm}

    {\normalsize
    \begin{tabular}{rl}
        \textbf{Autor:}          & Erlind Sejdiu \\
        \textbf{Matrikelnummer:} & 106357 \\[0.5cm]
    \end{tabular}
    }

    \vfill
    {\small\color{hkagray} Stand: \today}
\end{titlepage}

\newpage
\hypersetup{pageanchor=true}
\pagenumbering{arabic}
\setcounter{page}{1}

% =============================================================================
\section{Was ist Tauschhandel?}

Tauschhandel ist der direkte Austausch von Waren oder Leistungen ohne Geld.
Ware wechselt gegen Ware, Leistung gegen Leistung. Es gibt keinen festen
Preis, der Wert wird bei jedem Tausch neu ausgehandelt.

Smith~\cite{smith1776} und Menger~\cite{menger1892} beschreiben den Tausch als
Vorstufe des Geldes. Ein Tausch kommt nur zustande, wenn beide Seiten genau
das besitzen, was die andere Seite will. Dieses Problem heißt doppelte
Koinzidenz der Bedürfnisse. Als Lösung entstand nach diesem Modell Warengeld,
also Güter wie Salz, Getreide oder Vieh, die fast jeder annimmt.

Humphrey~\cite{humphrey1985} stellt fest, dass keine Gesellschaft belegt ist,
die allein auf Tauschhandel beruhte. Graeber~\cite{graeber2011} zeigt, dass
Gemeinschaften ihre Versorgung über Geschenke und Schulden regelten.
Tauschhandel ist vor allem in zwei Fällen belegt: beim Handel mit Fremden und
nach dem Zusammenbruch einer Währung. Beispiele für den zweiten Fall:

\begin{itemize}
    \item Deutschland 1945 bis 1948: Die Reichsmark verlor ihren Wert,
          Zigaretten dienten als Ersatzwährung. Die Währungsreform vom
          20.~Juni 1948 führte die D-Mark ein.
    \item Argentinien 2002: Nach der Finanzkrise gab es rund 4.500
          Tauschclubs in allen 23~Provinzen. Etwa 2,5~Millionen Menschen
          nutzten sie zur Versorgung~\cite{ambito2002}.
\end{itemize}

% =============================================================================
\section{Soziale Auffassung (Zeitgeist)}

Bargeld verliert an Bedeutung. 2023 wurden 51~Prozent der Einkäufe in
Deutschland bar bezahlt, 2021 waren es 58~Prozent. Mobiles Bezahlen stieg im
selben Zeitraum von 2 auf 6~Prozent. Nach Umsatz liegt die Debitkarte mit
32~Prozent vor dem Bargeld mit 26~Prozent~\cite{bundesbank2024}.

Tausch und Weitergabe von Dingen treten heute in zwei Formen auf: im Alltag
aus Spar- oder Umweltgründen und als Vorsorge für Krisen.

Zahlen zum Alltag:
\begin{itemize}
    \item Rund ein Viertel der Deutschen kauft aus zweiter Hand, Hauptgrund
          ist das Sparen (Kantar, 2.004~Befragte, 2024)~\cite{idealo2024}.
    \item 2024 wurden in Deutschland über 6,5~Milliarden Euro für
          Secondhand-Artikel ausgegeben, 52~Prozent davon für
          Kleidung~\cite{cebr2025}.
    \item Weltweit gibt es über 4.000 aktive Repair-Cafés und ähnliche
          Gruppen mit rund 190.000 erfolgreichen Reparaturen pro
          Jahr~\cite{nltimes2024}. Der Bund fördert solche Initiativen seit
          Dezember 2024 mit bis zu 3.000~Euro~\cite{heise2024}.
\end{itemize}

Zahlen zur Krise:
\begin{itemize}
    \item Die Inflationsrate lag 2022 bei 6,9~Prozent, dem höchsten Wert seit
          der Wiedervereinigung, und 2023 bei 5,9~Prozent~\cite{statistikportal}.
    \item Steigende Lebenshaltungskosten sind 2025 zum 15.~Mal die größte
          Sorge der Deutschen (52~Prozent von 2.400~Befragten)~\cite{ruv2025}.
    \item Im Oktober 2022 hatten 38~Prozent Lebensmittelvorräte angelegt,
          17~Prozent hatten keine Vorsorge getroffen~\cite{lmpraxis2022}.
          Das BBK empfiehlt einen Vorrat für zehn Tage. In Stuttgart
          erfüllen 29~Prozent diese Empfehlung~\cite{stuttgart2024}.
    \item Wegen der Sanktionen tauschen Russland und China wieder Waren
          direkt. Reuters zählte 2025 mindestens acht solcher Geschäfte,
          etwa chinesische Autos gegen russischen Weizen~\cite{reuters2025}.
\end{itemize}

% =============================================================================
\begin{landscape}

\newpage
\section{Themenkarte}

Die Karte zeigt, welche Begriffe zum Thema gehören und wie sie
zusammenhängen. Die Farben ordnen die Begriffe vier Bereichen zu. Jeder Pfeil
lässt sich als Satz lesen, vom Startbegriff über die Beschriftung zum
Zielbegriff.

\begin{figure}[H]
\centering
\begin{tikzpicture}[x=0.98cm, y=0.9cm]
    % --- Begriffe ---
    \node[kbox=clA] (A3) at (0,10)     {Misstrauen\\ in Geld};
    \node[kbox=clA] (A1) at (5.5,10)   {Weniger\\ konsumieren};
    \node[kbox=clC] (C1) at (11,10)    {Lebenswichtiges\\ (Essen, Medizin)};
    \node[kbox=clA] (A2) at (16.5,10)  {Für Krisen\\ vorsorgen};
    \node[kbox=clA] (A4) at (22,10)    {Wunsch nach\\ Gemeinschaft};

    \node[kbox=clB] (B1) at (5.5,6.5)  {Tauschpartner};
    \node[kbox=clB] (B2) at (11,6.5)   {Persönlicher\\ Wert};
    \node[kbox=clD] (D4) at (16.5,6.5) {Drohung\\ \& Gewalt};
    \node[kbox=clB] (B4) at (22,6.5)   {Vertrauen\\ \& guter Ruf};

    \node[kbox=clC] (C2) at (0,3)      {Genussmittel\\ (Tabak, Alkohol)};
    \node[kbox=clD] (D2) at (5.5,3)    {Sichere\\ Marktplätze};
    \node[kbox=clB] (B3) at (11,3)     {Feilschen};
    \node[kbox=clD] (D3) at (16.5,3)   {Weg \&\\ Transportrisiko};
    \node[kbox=clC] (C4) at (22,3)     {Erinnerungs-\\stücke};

    \node[kbox=clD] (D1) at (11,-0.5)  {Partnersuche\\ (Wer will was?)};
    \node[kbox=clC] (C3) at (16.5,-0.5){Arbeit \& Hilfe\\ (Reparieren, Schützen)};

    % --- Gesellschaft ---
    \draw[pfeil] (A1) -- node[txt]{schont} (C1);
    \draw[pfeil] (A1.north) |- ($(A4.north)+(0,0.75)$) node[txt, pos=0.75]{stärkt den} -- (A4.north);
    \draw[pfeil] (A2) -- node[txt]{hortet} (C1);
    \draw[pfeil] (A2) -- node[txt]{schützt vor} (D4);
    \draw[pfeil] (A3) -- node[txt, pos=0.6]{stärkt den} (B2);
    \draw[pfeil] (A3) -- node[txt, pos=0.45]{bringt zum\\ Tauschen} (B1);
    \draw[pfeil] (A4) -- node[txt]{baut auf} (B4);

    % --- Menschen ---
    \draw[pfeil] (B1) -- node[txt, pos=0.68]{beginnen mit\\ der} (D1);
    \draw[pfeil] (B1) -- node[txt, left=1pt, pos=0.25]{treffen sich auf} (D2);
    \draw[pfeil] ($(B2.south)+(-0.4,0)$) to[bend right=40] node[txt, anchor=east, pos=0.5]{setzt die\\ Grenze beim} ($(B3.north)+(-0.4,0)$);
    \draw[pfeil] ($(B3.north)+(0.4,0)$) to[bend right=40] node[txt, anchor=west, pos=0.5]{verändert\\ den} ($(B2.south)+(0.4,0)$);
    \draw[pfeil] (B4) -- node[txt, pos=0.4]{senkt} (D3);
    \draw[pfeil] (B4) -- node[txt]{verhindert} (D4);

    % --- Waren ---
    \draw[pfeil] (C1) -- node[txt]{erhöht bei\\ Mangel den} (B2);
    \draw[pfeil] (C2) -- node[txt, pos=0.4]{erhöhen durch\\ Sucht den} (B2);
    \draw[pfeil] (C2) -- node[txt, pos=0.4]{erleichtern als\\ Ersatzgeld die} (D1);
    \draw[pfeil] (C3) -- node[txt]{erleichtert\\ die} (D1);
    \draw[pfeil] (C4) -- node[txt]{stärken\\ als Pfand} (B4);

    % --- Hindernisse ---
    \draw[pfeil] (D1) -- node[txt, pos=0.55]{verlängert\\ den} (D3);
    \draw[pfeil] (D2) -- node[txt, pos=0.86]{verbieten} (D4);
    \draw[pfeil] (D3.south) |- (2.6,-2.0) node[txt, pos=0.75]{bremst die} |- (B1.west);
    \draw[pfeil] (D4) -- node[txt]{drückt\\ den} (B2);

    % --- Legende ---
    \foreach \c/\t [count=\i from 0] in {clA/Gesellschaft, clB/Menschen, clC/Waren, clD/Hindernisse} {
        \node[draw=black, line width=1pt, fill=\c!12, minimum width=0.5cm, minimum height=0.3cm, inner sep=0pt] at (\i*4+6,-2.8) {};
        \node[anchor=west, font=\scriptsize] at (\i*4+6.4,-2.8) {\t};
    }
\end{tikzpicture}
\caption{Themenkarte zum Thema Tauschhandel. Beispiel für das Lesen eines
Pfeils: \enquote{Lebenswichtiges erhöht bei Mangel den persönlichen Wert}.}
\end{figure}
\end{landscape}

% =============================================================================
\begin{thebibliography}{99}
\small
\bibitem{smith1776} Smith, A. (1776): \emph{An Inquiry into the Nature and
    Causes of the Wealth of Nations}. London: W.~Strahan \& T.~Cadell.
\bibitem{menger1892} Menger, C. (1892): On the Origin of Money. In:
    \emph{The Economic Journal}, 2(6), S.~239--255.
\bibitem{humphrey1985} Humphrey, C. (1985): Barter and Economic
    Disintegration. In: \emph{Man}, New Series, 20(1), S.~48--72.
\bibitem{graeber2011} Graeber, D. (2011): \emph{Debt: The First 5000 Years}.
    New York: Melville House.
\bibitem{ambito2002} Ámbito (2002): El trueque mueve ya 1.200 millones
    anuales en el país. Artikel vom 15.02.2002.
    \url{https://www.ambito.com/portada-principal/el-trueque-mueve-ya-1200-millones-anuales-el-pais-n3159776}
\bibitem{bundesbank2024} Deutsche Bundesbank (2024): Zahlungsverhalten in
    Deutschland 2023. Pressemitteilung vom 23.07.2024.
    \url{https://www.bundesbank.de/de/presse/pressemitteilungen/zahlungsverhalten-in-deutschland-2023-934828}
\bibitem{idealo2024} FashionUnited (2024): Ein Viertel der Deutschen kauft
    Secondhand (idealo Sparreport 2024).
    \url{https://fashionunited.de/nachrichten/mode/ein-viertel-der-deutschen-kauft-secondhand/2024072457455/}
\bibitem{cebr2025} Amazon / Centre for Economics and Business Research (2025):
    Second-Hand boomt. Pressemitteilung, Juli 2025.
    \url{https://press.aboutamazon.com/de/2025/7/second-hand-boomt-der-online-markt-in-europa-betraegt-21-6-milliarden-euro-trend-zeigt-weiter-nach-oben}
\bibitem{nltimes2024} NL Times (2024): Repair Café in Amsterdam celebrates
    15th anniversary. \url{https://nltimes.nl/node/75879}
\bibitem{heise2024} heise online (2024): Recht auf Reparatur: Repair-Cafés
    können ab Dezember Förderung beantragen. \url{https://heise.de/-9988463}
\bibitem{statistikportal} Statistische Ämter des Bundes und der Länder:
    Inflation. \url{https://www.statistikportal.de/de/inflation}
\bibitem{ruv2025} R+V Versicherung (2025): Die Ängste der Deutschen 2025.
    Pressemitteilung vom 18.09.2025.
    \url{https://www.ruv.de/newsroom/themenspezial-die-aengste-der-deutschen/pressemitteilungen/2025-09-18-studie-aengste-der-deutschen}
\bibitem{lmpraxis2022} Lebensmittel Praxis (2022): Menschen sorgen für
    Notvorräte vor dem Winter. Bericht vom 24.10.2022.
    \url{https://lebensmittelpraxis.de/handel-aktuell/35199-katastrophenschutz-menschen-sorgen-fuer-notvorraete-vor-dem-winter.html}
\bibitem{stuttgart2024} Landeshauptstadt Stuttgart (2024): Bevölkerungsumfrage
    zur Krisenvorsorge.
    \url{https://www.domino1.stuttgart.de/web/komunis/komunissde.nsf/a9830f07fd8c0e7dc1257cb70033e711/c66c90de6d4dd9a2c1258b040022ec33/$FILE/cwt01_.PDF}
\bibitem{reuters2025} Kyiv Independent (2025): Russia brings back barter trade
    to bypass sanctions, Reuters reports.
    \url{https://kyivindependent.com/russia-brings-back-barter-trade-to-bypass-sanctions-reuters-reports/}
\end{thebibliography}

\noindent{\small Alle Internetquellen abgerufen am 29.10.2026.}

\end{document}
